SGIRI es el sistema de tickets de mi proyecto: el cliente cuenta la falla, el administrador asigna un técnico y el técnico cierra con fotos de lo que hizo. El front está en Angular, el servidor en Spring Boot y los datos en PostgreSQL. También usa dos servicios de internet: Globalping, que mide la red desde sondas reales, y Cloudinary, que guarda las fotos y las firmas. En la segunda diapositiva muestro una petición de verdad: cambio el estado de un ticket con \texttt{PUT} y se ve el código que devuelve y lo que trae. En la tercera cuento qué pasa cuando Globalping no responde: el sistema intenta primero desde Ecuador por IPv4, después por IPv6 y después desde Sudamérica; si nada funciona, guarda el intento como fallido y repite el último dato bueno. Lo que saqué de esto es que cuando dependes de un servicio ajeno tienes que decidir qué mostrar si falla, y que dejar el fallo anotado es mejor que esconderlo.
